
\section*{2. Generalities on Formal Groups}
\phantomsection
\label{sga3:VIIB:2}
\addcontentsline{toc}{section}{2. Generalities on Formal Groups}

\subsection*{2.1}
\phantomsection
\label{sga3:VIIB:2.1}
\addcontentsline{toc}{subsection}{2.1. Formal groups and cogroups}

Let \(k\) be a pseudocompact ring and let \(G\) be a formal
\(k\)-group, that is, a group object in the category
\(\operatorname{Vaf}/k\) of formal varieties over \(k\). Let \(A\)
be the affine algebra of \(G\). The composition law of \(G\) plainly
defines a diagonal morphism, namely a homomorphism of profinite
\(k\)-algebras
\[
  \Delta_A:A\longrightarrow A\widehat\otimes_kA.
\]
This homomorphism has the following properties.

\begin{enumerate}
\item[\textup{\SGAInlineRefTarget{sga3:VIIB:2.1:item:i}(i)}] The diagram
\begin{center}
  \phantomsection
  \label{sga3:VIIB:2.1:coassociativity-square}
  \begin{tikzcd}[row sep=3em,column sep=4.7em]
A \arrow[r,"\Delta_A"] \arrow[d,"\Delta_A"']
  & A\widehat\otimes_k A
      \arrow[d,"{\Delta_A\widehat\otimes\operatorname{id}_A}"] \\
A\widehat\otimes_k A
  \arrow[r,"{\operatorname{id}_A\widehat\otimes\Delta_A}"]
  & A\widehat\otimes_k A\widehat\otimes_k A
\end{tikzcd}

\end{center}
is commutative.

\item[\textup{(ii)}] There is a necessarily unique
\emph{augmentation}, that is, a homomorphism of profinite
\(k\)-algebras
\[
  \varepsilon_A:A\longrightarrow k,
\]
such that the two composites displayed below are the identity maps of
\(A\):
\begin{center}
  \phantomsection
  \label{sga3:VIIB:2.1:counit-compositions}
  \(
\begin{gathered}
A\xrightarrow{\ \Delta_A\ }A\widehat\otimes_k A
 \xrightarrow{\ \varepsilon_A\widehat\otimes\operatorname{id}_A\ }
 k\widehat\otimes_k A\simeq A
\\[.7ex]
\text{and}\qquad
A\xrightarrow{\ \Delta_A\ }A\widehat\otimes_k A
 \xrightarrow{\ \operatorname{id}_A\widehat\otimes\varepsilon_A\ }
 A\widehat\otimes_k k\simeq A
\end{gathered}
\)

\end{center}

\item[\textup{\SGAInlineRefTarget{sga3:VIIB:2.1:item:iii}(iii)}] There is a necessarily unique \emph{antipode},
that is, a homomorphism of profinite \(k\)-algebras
\[
  c_A:A\longrightarrow A,
\]
such that the composite
\begin{center}
  \phantomsection
  \label{sga3:VIIB:2.1:antipode-composition}
  \(
A\xrightarrow{\ \Delta_A\ }A\widehat\otimes_k A
\xrightarrow{\ c_A\widehat\otimes\operatorname{id}_A\ }
A\widehat\otimes_k A
\xrightarrow{\ m_A\ }A
\)

\end{center}
equals \(\eta_A\circ\varepsilon_A\). Here \(m_A\) is the linear map
that sends \(a\widehat\otimes b\) to \(ab\), and
\(\eta_A:k\to A\) sends \(\lambda\) to \(\lambda1_A\).
\end{enumerate}

Conversely, data
\[
  (\Delta_A,\varepsilon_A,c_A)
\]
satisfying \textup{\SGARefLink{sga3:VIIB:2.1:item:i}{(i)}}--\textup{\SGARefLink{sga3:VIIB:2.1:item:iii}{(iii)}} endow \(G\) with the
structure of a formal \(k\)-group.\footnote{
\label{sga3:VIIB:nde:76}\textit{Editors' note.} One also says that \(A\) is a \emph{cogroup}
in the category of profinite \(k\)-algebras. The editors expanded the
material that follows, in particular by making explicit what a
morphism of formal groups \(K\to G\) means; cf.\ \SGARefLink{sga3:VIIB:2.3.1}{Proposition~2.3.1}.}
Explicitly, for every profinite \(k\)-algebra \(B\), the set
\[
  \operatorname{Hom}_c(A,B)
\]
of continuous morphisms of \(k\)-modules \(\phi:A\to B\) carries a
group structure, functorial in \(B\), defined by
\[
  \phi\cdot\phi'
  =
  m_B\circ
  (\phi\widehat\otimes\phi')
  \circ\Delta_A.
\]
The identity element is \(\eta_B\circ\varepsilon_A\), where \(m_B\)
is the multiplication of \(B\) and
\(\eta_B:k\to B\) sends \(\lambda\) to \(\lambda1_B\); the inverse
of \(\phi\) is \(\phi\circ c_A\). The set
\[
  \operatorname{Hom}_{\operatorname{Alp}/k}(A,B)
\]
of continuous morphisms of \(k\)-algebras \(A\to B\) is a subgroup,
because \(B\) is commutative.

\medskip
\noindent\textbf{Definition.}
\phantomsection
\label{sga3:VIIB:2.1.definition}
A \emph{morphism of formal \(k\)-groups}
\(\theta:K\to G\) is, by definition, a morphism of formal
\(k\)-varieties that respects the group structures. Let \(B\),
respectively \(A\), be the affine algebra of \(K\), respectively
\(G\), and let \(f:A\to B\) be the morphism corresponding to
\(\theta\). The condition is equivalently that \(f\) be compatible
with comultiplication:
\[
  (f\widehat\otimes f)\circ\Delta_A
  =
  \Delta_B\circ f.
\]
The identities
\[
  \varepsilon_B\circ f=\varepsilon_A,
  \qquad
  c_B\circ f=f\circ c_A
\]
then follow automatically. We write \(\operatorname{Grf}/k\) for
the category of formal \(k\)-groups.

\medskip
\noindent\textbf{Notation.}
\phantomsection
\label{sga3:VIIB:2.1.notation}
In what follows, the ideal
\[
  I_A=\operatorname{Ker}(\varepsilon_A)
\]
will be called the \emph{augmentation ideal} of \(A\), and we write
\[
  \omega_{G/k}=I_A/\overline{I_A^2}
\]
for the pseudocompact \(k\)-module obtained by quotienting \(I_A\)
by the closed ideal generated by the products \(xy\), with
\(x,y\in I_A\).
